% ================================================================ conclusions (v3)
\section{Conclusions}\label{sec:conclusions}

Specializing an open, general-purpose decision model on public radiology data improved its task-averaged accuracy on a radiology benchmark of report reading, case diagnosis and classification and radiological knowledge, and the specialized 27B model was more accurate than a medically trained LLM of the same size and than its own backbone used as an LLM. % [C1] [C4] [C5] [C16]
The effect differed between decisions: accuracy increased most in the detection of errors in reports and in Eurorad case diagnosis, a gain that derived from the answer options rather than from the case, whereas in the selection of the ACR Appropriateness Criteria panel, a catch-all option that was the most frequent training answer became the specialized models' default until their answer distributions were corrected for the training answer frequencies; averaged over tasks, specialization was as effective as a threefold increase in model size. % [C3] [C8] [C18]
The specialized models could answer more questions within a 5\% error budget than their starting points and the LLMs, but the larger model made more confident errors than the model it was specialized from. % [C15] [C12]
For categorical radiology decisions that must run within an institution, specializing a general-purpose decision model may be a practical alternative to scaling it up or to querying an LLM, provided that the answer spaces used for training and deployment are curated, including their catch-all options, and that the model is validated on the intended tasks and on local data. On an external test set of CT and MRI reports from another institution, specialization lowered the accuracy of the 27B model slightly and raised that of the 9B model. % [C22]
